\documentclass[12pt]{article}

\usepackage[a4paper, total={6in, 8in}]{geometry}
\usepackage{textcomp}

\begin{document}
\setlength{\parindent}{0pt}

\def \t {\textrightarrow}
\def \v {\vspace{1em}}
\def \bi {\begin{itemize}}
\def \ei {\end{itemize}}
\def \s[#1] {\section*{#1}}
\def \ss[#1] {\subsection*{#1}}
\def \sss[#1] {\subsubsection*{#1}}

\s[La pioggia nel pineto - D'Annunzio]
\ss[Versi 65-96]
Si crea una corrispondenza con il ripetersi dei verbi della sfera uditiva \t ascolta, odi \t corrispondence Boudelaire \\
Aeree cicale \t volano \t il piano è lo scroscio della pioggia, che mescolata a un canto delle cicale da lontano sembra salire \\
Si ripetono le dinamiche del chiasmo versi 75/76 \\
Le note \t sinestesia \t se si pensa all'aspetto sonoro, crea quasi una onomatopea \\
Non si sente la voce del mare \t che riporta al mare come oggetto, l'oggettistica dei poeti \t come nel naufragar me dolce di Leopardi, e anche in Pascoli \t il mare diventera ancora + oggetto in Montale \\
Ungaretti fa l'epurazione della parola d'annunziana \t la parola diventa una esternazione, come se nascesse in quel momento \t d'annunzio l'aveva resa impura, aveva sovraccaricata la parola con troppi sovrassensi \\
Evocazioni sonore sono evidenti \t pioggia d'argento ricorda Spleen \\
Piu folta men folta \t è un chiasmo \\
Chi sa dove! ripete la dinamica dell'uso della parola estremizzato \t anafora evidente, interrogativa indiretta \\
Limo \t è il fango che riporta alla belletta \\
La rana canta? \t no, gracchiano \t e le rana non ci sono d'estate \\
Limo rappresenta molti aspetti \t come quelli del fango (belletta), ma il limo è anche il fango del momento della creazione \t dio ha creato l'uomo con il fango \t lettura sacarale di d'annunzio \\

\ss[97-128]
Ciglia nere \t figura femminile estremizzata al sensualismo \\
Virente \t latinismo \t e anche aulente \\
Perche gli occhi come polle? \t perche saltano e scorrazzano \t e gli occhi sono curiosi di partecipare all'universo virente, che diventano parte di quel disegno naturale = panismo d'annunziano e naturalizzazione dell'uomo \\
Naturalizzazione priva forse della bellezza femminile \t denti come mandorle \\
Di fratta in fratta \t eco pascoliano, fru fru fra le fratte \t anche qua ripetizione della R \t che è anche una anafora \\
Intrica le ginocchia \t evoca il mondo del purgatorio, con personaggi come Manfredi che volevano sfuggire ma si ferirono cosi facendo \\
L'Opera si chiude con le stesse modalita con cui si è aperta \t visione concertistica di una sinfonia naturale, che coinvolge tutti i sensi \t e una complessità compositiva \\
Chiasmo continuo
\end{document}
